\section{Symmetry Selection Across Material Families}
  \label{sec:symmetry-selection-families}

  \subsection{General Symmetry Selection Principle}
    \label{subsec:general-selection}

    Within the projective framework,
    the superconducting gap symmetry is not postulated.
    It emerges from a minimization problem on the fiber
    raccordement cost under lattice and frustration constraints.

    Given a lattice symmetry group $G$
    and a dominant frustration channel characterized by
    a wavevector $\mathbf{Q}$,
    the admissible $w=2$ composite configurations
    belong to irreducible representations of $G$.
    The selected representation is the one that minimizes
    the total fiber raccordement cost functional,
    \begin{equation}
      \mathcal{C}_{\text{tot}}
      =
      \int_{\text{FS}}
      \mathcal{C}_{\Pi}
      \big(
      \mathbf{k};
      \mathbf{Q},
      \mathcal{F}
      \big)
      \, d\Omega ,
      \label{eq:cost-functional}
    \end{equation}
    subject to global fiber continuity.
    The explicit variational derivation of this minimization
    is given in Appendix~\ref{app:variational}.

    The functional $\mathcal{C}_{\Pi}$
    encodes the local mismatch between
    the composite phase configuration
    and the staggered frustration background.
    It depends on the dominant frustration vector $\mathbf{Q}$
    and on the frustration amplitude proxy $\mathcal{F}$.

    The symmetry channel minimizing Eq.~\ref{eq:cost-functional}
    defines the superconducting gap representation.

  \subsection{Role of the Frustration Ratio}
    \label{subsec:role-rF}

    To compare material families,
    we introduce the dimensionless frustration ratio
    \begin{equation}
      r_{\mathcal F}
      \equiv
      \text{first-order proxy for relative frustration amplitude},
      \label{eq:rF-definition}
    \end{equation}
    taken to scale proportionally,
    at leading order,
    with the magnetic structure factor $S(\mathbf{Q})$.

    The ratio $r_{\mathcal F}$ does not fix the symmetry by itself.
    It controls the relative weight of the staggered constraint
    in the cost functional $\mathcal{C}_{\Pi}$.

    Two limiting regimes can be identified:

    \begin{itemize}
      \item
      \textbf{Strongly staggered regime}
      ($r_{\mathcal F}$ large):
      the composite must strongly suppress overlap
      with the dominant frustration vector.
      Nodal structures aligned with $\mathbf{Q}$
      become energetically favorable.

      \item
      \textbf{Weak or isotropized frustration regime}
      ($r_{\mathcal F}$ small):
      the staggered constraint is reduced,
      and isotropic or extended $s$-wave
      configurations minimize the global cost.
    \end{itemize}

    In both cases,
    the lattice symmetry group $G$
    restricts the admissible irreducible representations.
    The interplay between $G$ and $r_{\mathcal F}$
    determines the selected gap symmetry.

  \subsection{Material Comparison Criterion}
    \label{subsec:material-comparison}

    Consider two material families
    with comparable lattice symmetry
    but different frustration amplitudes.
    Let their dominant exchange scales be
    $J_1$ and $J_2$,
    and their frustration ratios
    $r_{\mathcal F}^{(1)}$ and $r_{\mathcal F}^{(2)}$.

    To leading order,
    the amplitude scale of composite formation
    is set by $J$,
    while the symmetry channel
    is determined by the minimization of
    $\mathcal{C}_{\Pi}$
    under the corresponding $r_{\mathcal F}$.

    The transition temperature then scales as
    \begin{equation}
      T_c
      \sim
      \frac{\pi}{2}
      \delta J
      \cdot
      f(r_{\mathcal F}),
      \label{eq:Tc-general}
    \end{equation}
    where $f(r_{\mathcal F})$
    is a monotonic function encoding
    the effect of frustration amplitude
    on phase stiffness and composite stability.

    The precise functional form of $f$
    is model-dependent.
    However,
    its qualitative behavior is constrained:

    \begin{itemize}
      \item
      $f(r_{\mathcal F})$ increases
      from zero at vanishing frustration
      to a maximum at intermediate $r_{\mathcal F}$.

      \item
      For very large $r_{\mathcal F}$,
      phase stiffness is suppressed by reduced mobility,
      and $T_c$ decreases.
    \end{itemize}

    Equation~\ref{eq:Tc-general}
    provides a structural comparison tool
    between related compounds.

  \subsection{Preparation for Nickelates}
    \label{subsec:preparation-nickelates}

    Cuprates correspond to a regime
    of strong $(\pi,\pi)$ frustration
    and large $J$,
    yielding $d_{x^2-y^2}$ symmetry
    as the lowest-cost representation.

    Nickelates share similar lattice symmetry
    but exhibit reduced antiferromagnetic correlations,
    modified Fermi surface topology,
    and a smaller exchange scale.

    Within the present framework,
    these differences imply:

    \begin{enumerate}
      \item
      a reduced $r_{\mathcal F}$ relative to cuprates,

      \item
      a modification of the minimization landscape
      of $\mathcal{C}_{\Pi}$,

      \item
      a shift in the preferred irreducible representation.
    \end{enumerate}

    The next section develops
    the resulting quantitative and qualitative predictions.
